% !PRINTED 685
\weberpage{685}
%<<CONT
können wir die linearen Gleichungen, analog wie im \S.\,186, (5)
bilden:
\weq{2}{
\begin{aligned}
\xi_1 l_{1,1} + \xi_2 l_{2,1} + \cdots + \xi_{\nu-1} l_{\nu-1,1} + \delta_1 \xi_\nu &= \lambda_1, \\
\xi_1 l_{1,2} + \xi_2 l_{2,2} + \cdots + \xi_{\nu-1} l_{\nu-1,2} + \delta_2 \xi_\nu &= \lambda_2, \\
\cdot\ \cdot\ \cdot\ \cdot\ \cdot\ \cdot\ \cdot\ \cdot\ \cdot\ \cdot\ \cdot\ \cdot & \\
\xi_1 l_{1,\nu} + \xi_2 l_{2,\nu} + \cdots + \xi_{\nu-1} l_{\nu-1,\nu} + \delta_\nu \xi_\nu &= \lambda_\nu,
\end{aligned}
}
dessen Determinante [\S.\,186, (4)] den von Null verschiedenen
Werth hat:
\weq{3}{
\left|
\begin{array}{ccccc}
l_{1,1}, & l_{2,1}, & \ldots, & l_{\nu-1,1}, & \delta_1\\
l_{1,2}, & l_{2,2}, & \ldots, & l_{\nu-1,2}, & \delta_2\\
\cdot & \cdot & \cdot & \cdot & \cdot\\
l_{1,\nu}, & l_{2,\nu}, & \ldots, & l_{\nu-1,\nu}, & \delta_\nu
\end{array}
\right| = n\,L\,(\varepsilon_1, \varepsilon_2, \ldots, \varepsilon_{\nu-1})
}
Demnach sind die Unbekannten $\xi_1, \xi_2, \ldots, \xi_{\nu-1}, \xi_\nu$ aus (2)
\gesperrt{eindeutig} bestimmt.

Durch Addition der Gleichungen (2) findet man zunächst
sehr einfach
\weq{4}{
n\,\xi_\nu = \log N_a(\alpha),
}
und $\xi_\nu$ bleibt also dasselbe für alle \gesperrt{Zahlen $\alpha$, die unter einander
 associirt sind.}

Die Zahlen $\xi_1, \xi_2, \ldots, \xi_{\nu-1}$ heissen die Exponenten der
Zahl $\alpha$ [in Bezug auf das Fundamentalsystem $(\varepsilon_1, \varepsilon_2, \ldots, \varepsilon_{\nu-1})$].

Aus dieser Definition ergiebt sich ohne Weiteres, dass die
Exponenten eines Productes oder eines Quotienten gleich der
Summe oder der Differenz der entsprechenden Exponenten der
einzelnen Bestandtheile sind. Die Exponenten einer Einheit sind
ganze rationale Zahlen, und jedes System ganzer rationaler Zahlen
$\xi_1, \xi_2, \ldots, \xi_{\nu-1}$ ist auch das Exponentensystem einer Einheit.

Die Exponenten associirter Zahlen unterscheiden sich also
um ganze Zahlen von einander, und man kann zu jeder Zahl $\alpha$
eine associirte Zahl $\alpha_0$ finden, deren Exponenten zwischen 0
und 1 liegen. Solche Zahlen heissen \gesperrt{reducirte Zahlen} (in
Bezug auf das Fundamentalsystem $\varepsilon_1, \ldots, \varepsilon_{\nu-1}$). Reducirte Einheiten
sind alle und nur die in $\Omega$ vorhandenen \gesperrt{Einheitswurzeln,
deren Zahl wir mit $w$ bezeichnen wollen.} Da
die beiden Einheitswurzeln $\pm 1$ in jedem Körper enthalten sind,
so ist $w$ mindestens $= 2$ und immer eine gerade Zahl.

Wenn zwei associirte Zahlen dasselbe Exponentensystem
haben, so unterscheiden sie sich nur durch einen Factor, der
%CONT>>
